\section{Hallucination Has a Formal Definition}
\label{sec:hallucination}

Hallucination is typically described informally: the model says something
false or unsupported. This is a behavioral description, not a mechanistic
one. It gives no account of why hallucination happens or what would prevent
it.

\subsection*{The BP Definition}

The BP framework provides a formal definition: a model hallucinates when
its output beliefs diverge from the correct posterior given the evidence.
For a grounded transformer with a declared ontology and verifier, this
divergence is measurable:
\[
b(j) \neq P_{\mathrm{true}}(j),
\]
where $P_{\mathrm{true}}$ is the correct posterior induced by the knowledge
base and observed evidence.

For an ungrounded transformer, the comparison is not well-defined because
there is no formally declared ground truth. Hallucination in ungrounded
models is not a numerical error in the inference procedure --- it is a
consequence of operating without a fully specified and verifiable grounding
structure.

\subsection*{What Grounding Enables --- and Its Limits}

A fully grounded transformer would need: a declared concept space, a
declared relational structure, a grounding procedure mapping natural
language to the declared system, and a verification mechanism.

Grounding does not eliminate all possible error. Even a formally grounded
system can have an incomplete ontology, incorrect premises, ambiguous
mappings from natural language, or approximation errors in inference.
What grounding provides is not a guarantee of truth but a guarantee of
\emph{verifiable correctness relative to a declared ontology}: outputs can
be checked against the knowledge base, and failures can be traced to
specific premises or inference steps. This is a weaker but more honest
claim --- and still substantially stronger than what ungrounded systems
can provide.

Retrieval-augmented generation is a partial form of grounding: it
temporarily constrains the inference space but does not provide persistent
ontological structure or a verifier.

\subsection*{Epistemic Status}

\begin{center}
\begin{tabular}{ll}
\toprule
Claim & Status \\
\midrule
Grounded BP permits formally verifiable correctness & Direct consequence \\
Grounding eliminates all hallucination & Explicitly not claimed \\
Ungrounded models lack a formal correctness criterion & Direct consequence \\
Hallucination requires grounding to formally address & Mechanistic interpretation \\
Scaling alone cannot eliminate hallucination & Mechanistic interpretation \\
\bottomrule
\end{tabular}
\end{center}